\begin{proofbox}
\textbf{\textcolor{CProof}{证明一：最快证明（考场 2--3 步，首选）}}
\begin{center}
\begin{tikzpicture}[font=\small]
 \path[use as bounding box] (-4,-2.9) rectangle (4,3);
 \node at (0,2.7) {两条连线共用高度 $y$，水平跨度分别为 $x+a$ 与 $a-x$};
 \draw[blue!60!black,thick] (0,0) ellipse (3 and 2);
 \coordinate (A) at (-3,0);\coordinate (B) at (3,0);\coordinate (P) at (1.8,1.6);\coordinate (H) at (1.8,0);
 \draw[teal!70!black,very thick] (A)--(P);
 \draw[magenta!70!black,very thick] (B)--(P);
 \draw[gray,thick] (A)--(B);
 \draw[dashed,thick] (P)--(H) node[midway,left] {$y$};
 \draw (H) ++(-.18,0)--++(0,.18)--++(.18,0);
 \node[below left] at (A) {$A_1$};\node[below right] at (B) {$A_2$};
 \node[above right] at (P) {$P(x,y)$};\node[below] at (H) {$H$};
 \node[above] at (-.7,0) {$x+a$};\node at (2,-.8) {$a-x$};
 \node at (0,-2.55) {$k_1k_2=-\frac{y^2}{(x+a)(a-x)}=-\frac{b^2}{a^2}$};
\end{tikzpicture}
\end{center}


\par\textbf{切入信号：}题目要求两条连线的\textbf{斜率乘积}，先乘再化，不分别求具体斜率。
设 $P(x,y)\in C\setminus\{A_1,A_2\}$，则
\[
 \begin{aligned}
 k_{PA_1}k_{PA_2}
 &=\frac{y}{x+a}\cdot\frac{y}{x-a}
 =\frac{y^2}{x^2-a^2},\\
 y^2&=b^2\Bigl(1-\frac{x^2}{a^2}\Bigr)
 =-\frac{b^2}{a^2}(x^2-a^2),\\
 \therefore\quad k_{PA_1}k_{PA_2}&=\boxed{-\frac{b^2}{a^2}}.
 \end{aligned}
\]
\textbf{最关键的一眼：}$y^2$ 与 $x^2-a^2$ 在椭圆方程中\textbf{恰好成比例}，不必再设第二个未知量。

\end{proofbox}
\clearpage
\begin{proofbox}
\textbf{\textcolor{CProof}{证明二：最低门槛（圆的直角 $\longrightarrow$ 椭圆的斜率积）}}
\begin{center}
\begin{tikzpicture}[font=\small]
 \path[use as bounding box] (-2.8,-3.9) rectangle (9.8,3.3);
 \foreach \shift/\ry/\name in {0/2/圆,7/1.2/椭圆}{
 \begin{scope}[shift={(\shift,0)}]
 \node at (0,2.95) {\name：横坐标保持不变};
 \draw[->,gray] (-2.4,0)--(2.4,0);
 \draw[blue!60!black,thick] (0,0) ellipse (2 and \ry);
 \coordinate (L) at (-2,0);\coordinate (R) at (2,0);\coordinate (T) at (1.2,{.8*\ry});
 \draw[teal!70!black,thick] (L)--(T);
 \draw[magenta!70!black,thick] (R)--(T);
 \draw[dashed] (1.2,0)--(T);
 \fill (T) circle (1.5pt);
 \node[below left] at (L) {$A_1$};\node[below right] at (R) {$A_2$};
 \end{scope}}
 % Explicit right angle on the circle, with its own coordinates.
 \coordinate (CL) at (-2,0);\coordinate (CR) at (2,0);\coordinate (Q) at (1.2,1.6);
 \pic[draw,angle radius=3mm] {right angle=CL--Q--CR};
 \node[above right] at (Q) {$Q(u,v)$};
 \node[above right] at (8.2,.96) {$P(u,\frac ba v)$};
 \node at (0,-2.6) {$k_{QA_1}k_{QA_2}=-1$};
 \node at (7,-2.6) {$k_{PA_1}k_{PA_2}=-\bigl(\frac ba\bigr)^2$};
 \draw[->,thick] (2.65,1.95)--(4.35,1.95) node[midway,above] {$y\mapsto\frac ba y$};
 \node at (3.5,-3.4) {图取 $a=2,b=1.2$，两侧单位长度相同};
\end{tikzpicture}
\end{center}


\textbf{要回答的三个问题：}点是怎样配对的？为什么比例恰是 $b/a$？缩放后究竟保留了什么？

\begin{description}[style=nextline,leftmargin=2.8em,labelsep=0.5em,itemsep=0.55em]
 \item[\textbf{第 1 步：搭一座共同的桥——固定两端点。}]
  取圆 $C_0:u^2+v^2=a^2$，以及椭圆
  $C:\frac{x^2}{a^2}+\frac{y^2}{b^2}=1\ (a>b>0)$。
  它们的中心都是 $O(0,0)$，左右端点同为 $A_1(-a,0),A_2(a,0)$。
  选圆上点 $Q(u,v)\ne A_1,A_2$；于是 $v\ne0$，且 $-a<u<a$。

 \item[\textbf{第 2 步：先在圆里得到一个熟悉的 $-1$。}]
  $A_1A_2$ 是直径，由圆周角定理，$\angle A_1QA_2=90^\circ$。
  两条斜率存在，因而
  \[
    k_{QA_1}=\frac{v}{u+a},\qquad
    k_{QA_2}=\frac{v}{u-a},\qquad
    \boxed{k_{QA_1}k_{QA_2}=-1}.
  \]
  （也可直接算 $\frac{v^2}{u^2-a^2}=\frac{a^2-u^2}{u^2-a^2}=-1$。）

 \item[\textbf{第 3 步：为什么纵坐标要乘 $b/a$？不是随便猜的。}]
  我们\textbf{主动选择保持横坐标不变}，令圆上的 $Q(u,v)$ 对应到椭圆上的
  $P(u,w)$，并要求两点在横轴同侧。由于 $u^2+v^2=a^2$，而
  $P$ 要满足椭圆方程，故
  \[
    \frac{u^2}{a^2}+\frac{w^2}{b^2}=1
    \ \Longrightarrow\ 
    \frac{w^2}{b^2}=\frac{v^2}{a^2}
    \ \Longrightarrow\ 
    w=\frac ba\,v
    \quad\text{（同侧取同号）}.
  \]
  因此映射只能写成（在上述规则下）
  \[
    \boxed{T:\ (u,v)\mapsto\left(u,\frac ba v\right)}.
  \]
  \textbf{注意：}并非所有可能的圆椭圆对应都必须如此；我们选这种映射，是因为它\textbf{固定横轴}。

 \item[\textbf{第 4 步：识别不变量与变化量。}]
  横坐标相同，垂足 $H=(u,0)$ 不动；横轴所有点也固定：
  $T(A_1)=A_1,\ T(O)=O,\ T(A_2)=A_2$。
  $Q$ 到横轴的有向高度 $v$ 变成 $\frac ba v$，纵向增量都乘 $b/a$。
  归一化高度保持 $\frac{y_P}{b}=\frac{v_Q}{a}$。
  从端点到 $Q$ 的连线经 $T$ 后正好变成从同一端点到 $P$ 的连线。

 \item[\textbf{第 5 步：两条斜率为什么\emph{各自}乘一次 $b/a$？}]
  对左端点，水平位移仍是 $u+a$，竖直位移由 $v$ 变为 $(b/a)v$，所以
  \[
    k_{PA_1}=\frac{(b/a)v}{u+a}
    =\frac ba\,k_{QA_1}.
  \]
  对右端点同理（水平位移是 $u-a$，注意它是负数）：
  \[
    k_{PA_2}=\frac{(b/a)v}{u-a}
    =\frac ba\,k_{QA_2}.
  \]
  \textbf{这是核心几何意义：}水平跨度原封不动，高度等比变小，\textbf{斜率也就按比例变小}。

 \item[\textbf{第 6 步：斜率积随之缩放两次。}]
  \[
   \begin{aligned}
    k_{PA_1}k_{PA_2}
    &=\left(\frac ba\,k_{QA_1}\right)
      \left(\frac ba\,k_{QA_2}\right)\\
    &=\left(\frac ba\right)^2(-1)
      =\boxed{-\frac{b^2}{a^2}}.
   \end{aligned}
  \]
 \item[\textbf{第 7 步：逆向检验——椭圆也能拉回圆。}]
  反过来给定椭圆点 $P(x,y)$，令 $Q=(x,\frac ab y)$，
  则 $x^2+(\frac ab y)^2=a^2$，即 $Q$ 在圆上，且 $T(Q)=P$。
  因而这套对应对每个点都成立，是\textbf{可逆的一一对应}，并非只对图中某个特选点成立。
\end{description}

\textbf{两条必要提醒：}① 圆里 $QA_1\perp QA_2$，但压缩后的
$PA_1,PA_2$ \textbf{一般不再垂直}；本证明用的是斜率的\textbf{缩放关系}，
不是用“压缩后仍为直角”。② 排除端点，是因为端点处一条连线退化、斜率无定义。

\textbf{为什么推荐此法：}初中圆周角定理 $\to$ 统一的纵向缩放 $\to$ 高中斜率乘积，
不必死记公式，也能理解负号与平方的来源。
\end{proofbox}
\clearpage
\begin{proofbox}
\textbf{\textcolor{CProof}{证明三：常规解析法（步骤完整，适合书面证明）}}
\begin{center}
\begin{tikzpicture}[font=\small]
 \path[use as bounding box] (-4,-3.1) rectangle (4.5,3.1);
 \node at (0,2.8) {常规法：从两个端点分别写出坐标差};
 \draw[->,gray] (-3.6,0)--(3.7,0) node[right] {$x$};
 \draw[->,gray] (0,-2.25)--(0,2.25) node[left] {$y$};
 \draw[blue!60!black,thick] (0,0) ellipse (3 and 2);
 \draw[teal!70!black,very thick] (-3,0)--(-1.8,-1.6);
 \draw[magenta!70!black,very thick] (3,0)--(-1.8,-1.6);
 \draw[dashed] (-1.8,0)--(-1.8,-1.6);
 \node[below left] at (-1.8,-1.6) {$P(x_0,y_0)$};
 \node[above left] at (-3,0) {$A_1(-a,0)$};\node[above right] at (3,0) {$A_2(a,0)$};
 \node[above] at (-1.8,0) {$x_0$};
 \node[teal!70!black] at (-3.5,-1.3) {$k_1<0$};
 \draw[teal!70!black,thin] (-3.05,-1.1)--(-2.5,{-.5*4/3});
 \node[magenta!70!black] at (1,-1.2) {$k_2>0$};
 \node at (0,-2.7) {$k_1=\frac{y_0}{x_0+a},\quad k_2=\frac{y_0}{x_0-a},\quad x_0\ne\pm a$};
\end{tikzpicture}
\end{center}


\begin{description}[style=nextline,leftmargin=2.8em,labelsep=0.5em,itemsep=0.4em]
 \item[\textbf{第 1 步（设点并检查定义域）：}]
  设 $P(x_0,y_0)$，$\frac{x_0^2}{a^2}+\frac{y_0^2}{b^2}=1$，$P\ne A_1,A_2$。
  所以 $x_0\ne\pm a$，两条连线斜率均存在。
 \item[\textbf{第 2 步（两点斜率公式）：}]
  \[
   k_{PA_1}=\frac{y_0}{x_0+a},\qquad k_{PA_2}=\frac{y_0}{x_0-a}.
  \]
 \item[\textbf{第 3 步（先乘，平方差结构）：}]
  \[
   k_{PA_1}k_{PA_2}=\frac{y_0^2}{(x_0+a)(x_0-a)}
   =\frac{y_0^2}{x_0^2-a^2}.
  \]
 \item[\textbf{第 4 步（椭圆方程消元）：}]
  \[
   y_0^2=b^2\Bigl(1-\frac{x_0^2}{a^2}\Bigr)
   =\frac{b^2}{a^2}(a^2-x_0^2).
  \]
 \item[\textbf{第 5 步（约分，判断符号）：}]
  \[
   k_{PA_1}k_{PA_2}
   =\frac{b^2}{a^2}\frac{a^2-x_0^2}{x_0^2-a^2}
   =-\frac{b^2}{a^2}.
  \]
  其中约分合法的原因是 $x_0\ne\pm a$。
\end{description}

\end{proofbox}
\clearpage
\begin{proofbox}
\textbf{\textcolor{CProof}{补充证明四：长轴在 $y$ 轴时直接检验（防止套错）}}
\begin{center}
\begin{tikzpicture}[font=\small]
 \path[use as bounding box] (-3,-4.2) rectangle (5,4.3);
 \node at (.6,3.9) {长轴竖直：两条斜率共用分母 $x$};
 \draw[->,gray] (-2.5,0)--(3,0) node[right] {$x$};
 \draw[->,gray] (0,-3.4)--(0,3.35) node[left] {$y$};
 \draw[blue!60!black,thick] (0,0) ellipse (2 and 3);
 \coordinate (L) at (0,-3);\coordinate (R) at (0,3);\coordinate (P) at (1.6,1.8);
 \draw[teal!70!black,very thick] (L)--(P);
 \draw[magenta!70!black,very thick] (R)--(P);
 \fill (L) circle (1.5pt) node[below left] {$B_1(0,-a)$};
 \fill (R) circle (1.5pt) node[above right] {$B_2(0,a)$};
 \fill (P) circle (1.5pt) node[right] {$P(x,y)$};
 \node[teal!70!black,right] at (2.3,-1.1) {$k_1=\frac{y+a}{x}$};
 \node[magenta!70!black,right] at (2.3,2.6) {$k_2=\frac{y-a}{x}$};
 \node at (.5,-3.85) {图取 $a=3,b=2$：$k_1k_2=3\cdot(-\frac34)=-\frac94$};
\end{tikzpicture}
\end{center}


对 $\dfrac{y^2}{a^2}+\dfrac{x^2}{b^2}=1$，
$B_1(0,-a),B_2(0,a)$，$P(x,y)\ne B_1,B_2$，故 $x\ne0$：
\[
 \begin{aligned}
 k_{PB_1}k_{PB_2}
 &=\frac{y+a}{x}\cdot\frac{y-a}{x}
 =\frac{y^2-a^2}{x^2}\\
 &=\frac{-\frac{a^2}{b^2}x^2}{x^2}
 =\boxed{-\frac{a^2}{b^2}}.
 \end{aligned}
\]
\textbf{作用：}可以从斜率的“分子/分母”重新推出公式，而不是硬记一个倒数。

\end{proofbox}
\clearpage
\begin{proofbox}
\textbf{\textcolor{CProof}{补充证明五：逆命题与轨迹方程（高价值拓展）}}
\begin{center}
\begin{tikzpicture}[font=\small]
 \path[use as bounding box] (-2.1,-3) rectangle (10.1,3);
 \node at (4,2.65) {固定水平两点：参数决定它们是长轴端点还是短轴端点};
 \foreach \sx/\ry/\la/\role in {0/.6/{\frac14}/长轴端点,4/1.2/{1}/直径端点,8/1.8/{\frac94}/短轴端点}{
 \begin{scope}[shift={(\sx,0)}]
 \draw[gray,->] (-1.6,0)--(1.6,0);
 \draw[blue!60!black,thick] (0,0) ellipse (1.2 and \ry);
 \draw[magenta!70!black,fill=BGProof] (-1.2,0) circle (2pt) (1.2,0) circle (2pt);
 \node at (0,-2.2) {$\lambda=\la$};
 \node at (0,-2.75) {\role};
 \end{scope}}
 \node at (4,1.95) {圆};
 \node at (4,-3.4) {同尺度；空心端点均须从轨迹中删除};
\end{tikzpicture}
\end{center}


固定 $A_1(-a,0),A_2(a,0)$，设斜率积为 $-\lambda$（$\lambda>0$）。
由斜率公式得
\[
 \frac{y^2}{x^2-a^2}=-\lambda
 \quad\Longleftrightarrow\quad
 \frac{x^2}{a^2}+\frac{y^2}{\lambda a^2}=1,
 \qquad x\ne\pm a.
\]
\textbf{结论：}斜率积为 $-\lambda$ 刻画了一条去掉指定端点的椭圆（$\lambda\ne1$）或圆（$\lambda=1$）；
$0<\lambda<1$ 才能判定 $A_1,A_2$ 是\textbf{长轴}端点。

\medskip
\textbf{\textcolor{CProof}{方法选择}}
\begin{itemize}[leftmargin=*,itemsep=0.2em]
 \item \textbf{赶时间写答案：}选最快法，直接写 $y^2=-\frac{b^2}{a^2}(x^2-a^2)$。
 \item \textbf{想真正记住：}选圆 $\to$ 椭圆压缩法，牢记斜率各缩小 $b/a$。
 \item \textbf{要求严谨分步：}选常规法，先核查 $P$ 不是长轴端点。
 \item \textbf{反求轨迹或参数：}选逆向法，先用乘积构造方程，再看 $\lambda$ 的范围。
\end{itemize}
\end{proofbox}
